\chapter{Introduction \& Background}

\ac{hiv1}, a virus transmitted by body fluids, capable of shutting down the host immune system if not treated immediately \cite{scutchfieldAIDSUpdate1989}. As of 2020, an estimated 37 million people around the globe suffer from \ac{hiv1} \cite{venturaHumanImmunodeficiencyVirus2020}. Furthermore, \ac{hiv1} continues to spread, where it is estimated an additional 1.7 million global citizens acquire \ac{hiv1} on average each year \cite{venturaHumanImmunodeficiencyVirus2020}. Despite having been discovered over 4 decades ago in 1981 \cite{sharpOriginsHIVAIDS2011}, \ac{hiv1}, the cause of \ac{aids}, remains an incurable disease and evades any form of vaccination or cure \cite{venturaHumanImmunodeficiencyVirus2020} where the best treatment option for patients of \ac{hiv1} is to undergo \ac{art} treatment \cite{phanuphakHIVTreatmentPrevention2020}. With \ac{art}, \ac{hiv1} associated mortality rates have been successfully reduced ultimately saving millions, and changing \ac{aids} from a deadly disease, to one where the patients may typically live a full life \cite{venturaHumanImmunodeficiencyVirus2020}. However, it is critical to accept \ac{art} not as a complete solution, rather, as a temporary solution. \ac{art} itself remains an imperfect treatment with problems ranging from rebound due to patients not fulfilling the required periodic treatment plan \cite{buellLifelongAntiretroviralTherapy2016}, to constrained availability of treatment \cite{schoutenAntiretroviralDrugSupply2011}, to problematic cost of treatment \cite{kahnCostEffectivenessAntiretroviralTherapy2011,sharpOriginsHIVAIDS2011}.

\section{Drug Resistance}

Amongst these problems, however, there underlies a biologically critical change \ac{hiv1} is undergoing common to many pathogens where a treatment is not potent enough to cure is the gold standard and thus deterring the progress of \ac{art}. Specifically, in recent years, with advent of \ac{art}, \ac{hiv1}, has been slowly but surely developing resistance to many modern \ac{art} treatment related drugs \cite{clutterHIV1DrugResistance2016,fergusonHIV1ReplicationCycle2002,HIV1GeneticVariation2013}.

To elaborate on \ac{hiv1} \acp{drm}, it is typical that a single amino acid variation coded for in the \ac{hiv1} genome is not enough to induce resistance towards \acp{ard} \cite{clutterHIV1DrugResistance2016}. Furthermore, within a given drug classification, it is rather common to see cross-resistance where the increased resistance towards one particular drug increases the resistance towards another \cite{clutterHIV1DrugResistance2016}. Despite such cross-resistance prevalence within the same drug class, the same cannot be said across different drug classes where current research seems to indicate that it is more difficult \cite{larderInteractionsDrugResistance1994} for such cross-resistances, however, not impossible \cite{clutterHIV1DrugResistance2016,lombardiPrevalenceFactorsAssociated2021,segal-maurerCapsidInhibitionLenacapavir2022}. This lower prevalence of cross-class drug resistance is a key reason for the use of multiple \acp{ard} of differing classes during modern \ac{art} treatment plans \cite{gibasTwodrugRegimensHIV2022}. Notably, such \acp{drm} seem to typically result in a fitness penalty to the host strain of the \acp{drm} \cite{clutterHIV1DrugResistance2016}.

\section{Immune Escape}

\Ac{hiv1} is fascinating in that it is one of the few viruses in modern day that both is highly prevalent in the population, but also eludes cures. Specifically, \ac{ctl}, also known as CD8+ T cells, are well known to be a host's main antagonist to the\ac{hiv1} virus \cite{garciaEffectInterferenceCD82015}. Then, it goes to reason, that \ac{hiv1}, with a relatively unstable genome capable of high rates of proliferation \cite{richmanHIVEvolutionEscape2004}, since it's initial entry into the human population, has developed various mechanisms of circumvent, or escape, an effective \ac{ctl} response. One such mechanism is via certain specific \acp{iem} that result in the suppressing of displaying specifically \ac{hiv1} associated antigens thus avoiding apoptosis via natural killer T cells \cite{varela-rohenaControlHIV1Immune2008}. Similar \acp{drm}, \acp{iem} such \acp{iem} also seem to introduce a fitness penalty to the \ac{hiv1} strain \cite{woodmanHIVMolecularEpidemiology2009,leslieHIVEvolutionCTL2004}.

\section{Motivation for Study}

Both mutations, independent of each other are relatively thoroughly studied. However, at the time of writing, there exists a minimal amount of studies looking into the intersection between both \acp{drm} and \acp{iem}. In other words, since both forms of mutations incur fitness penalties, an area of interest may be how \acp{drm} affect the mutation based immune escape process via epitopic modifications. The few studies emphasizing such a topic discusses various \acp{drm} that both corroborate the increase of immunogenicity of \ac{hiv1} \cite{masonAntiretroviralDrugResistance2004,blanco-herediaIdentificationImmunogenicCytotoxic2016}, and a decrease in immunogenicity, depending on varying factors \cite{muellerInfluenceMajorHIV12011,rahmanRaltegravirElvitegravirresistanceMutation2014}. The results of these studies are reasonable as \acp{drm} occurring within epitope regions, i.e., regions that affect the presentation via \ac{hla1} may modulate the fitness penalty due to directly affecting a immune escape mechanism. This study hopes to extend the knowledge base surrounding such interplay between \acp{drm} and \acp{iem} as such knowledge may serve foundational to improving treatment regiment by selecting drugs that produce strong fitness penalties and the future development of \ac{hiv1} vaccines, as explained by previous studies \cite{muellerInfluenceMajorHIV12011,rahmanRaltegravirElvitegravirresistanceMutation2014}.